\documentclass[10pt]{article}

\usepackage{amsmath}
%
\usepackage{amssymb}
%
\usepackage[width=7in,height=10in,top=0.75in,centering]{geometry}
%
\usepackage[dvipsnames]{xcolor}
%
\usepackage{hyperref}
%
%%%%%%%%%%%%%%%% HEADER %%%%%%%%%%%%%%%%%%%%%
\usepackage{fancyhdr}
\pagestyle{fancy}
%%
%\lfoot{Left footer}
%\cfoot{Center of footer}
%\rfoot{Right of footer}
%%
\lhead{Survey of Calculus I}
\chead{Knowledge Checkpoint 3 Reflection}
\rhead{\thepage}
%
\parindent0pt
\fboxsep=0.25cm
%
\newcommand{\Heading}[2]{\fbox{{\Large\textbf{#1}}}\hfill\fbox{\textbf{#2}}\par\vspace{0.1in}}
%
\newcommand{\Name}[0]{\textbf{Name}:\,\hrulefill\hrulefill\,\,}
%
\newcommand{\Date}[0]{\textbf{Date}:\,\hrulefill\,\,}
%
\newcommand{\Section}[0]{\textbf{Section}:\,\hrulefill\,\,}

%%%%%%%%%%%% BEGIN DOCUMENT %%%%%%%%%%%%%%%
%%%%%%%%%%%% BEGIN DOCUMENT %%%%%%%%%%%%%%%
%%%%%%%%%%%% BEGIN DOCUMENT %%%%%%%%%%%%%%%
%%%%%%%%%%%% BEGIN DOCUMENT %%%%%%%%%%%%%%%
%%%%%%%%%%%% BEGIN DOCUMENT %%%%%%%%%%%%%%%

\begin{document}

\thispagestyle{empty}

\Heading{MATH 1190/1210 Survey of Calculus I}{\begin{minipage}{0.25\textwidth}\begin{center} Knowledge Checkpoint 3 \\Reflection\end{center}\end{minipage}}

\vspace{0.1in}

\Name{}\Date{}\Section{}

\hrulefill
%
\paragraph*{GOAL:} Struggle and making mistakes are healthy parts of growth as long as both are framed productively. The intention of this activity is to create an opportunity for you to re-frame struggle \& mistakes productively so you can grow as a learner.\\

{\bf That you are honest with yourself when you respond to these prompts is a critically important for your growth.} This can be difficult, so please just do your best.\\

\hrulefill\\

\begin{enumerate}

\item Read through the Knowledge Checkpoint 3 Solutions on Canvas, and briefly review course materials. {\bf After doing so, which course experiences} (e.g. pre-class assignments, lessons, homework, practice problems) {\bf do you think would have best prepared someone for this knowledge checkpoint?}

    \vfill

\item The learning targets A2, A3, I1, and I2 will re-appear in Knowledge Checkpoint 4. For each of these learning targets, what mistakes (if any) did you make in Knowledge Checkpoint 3, and how could you avoid such a mistake next time?

\vfill

\vfill

Regardless, please know that {\it making mistakes is \underline{not} a sign of weakness} and {\it mistakes do not reflect poorly on you as a person}. {\bf Mistakes are just opportunities to  grow!} :)

% %%%%%%%%%
% \newpage%
% %%%%%%%%%

% \item On Tuesday 3 December at 7:00 PM, we will have assessment 4. Assessment 4 will cover targets {\bf\color{Fuchsia}I4}, {\bf\color{Orange}P1}, {\bf\color{Orange}P2}, and {\bf\color{Orange}P3}. 

% On Wednesday 11 December at 7:00 PM, we will have our final assessment. You can find a list of the learning targets that will appear on the final assessment by visiting the ``Learning Targets" page on Canvas. 

% {\bf Make a day-by-day plan to prepare for both assessments on the next two pages.} Include which target(s) you will study on each listed day, what resources you will use to study that target, and how you will use good study practices when engaging with those resources. 

% In exam wrappers and the ``How to Study" video series, we have seen several evidence-based study practices. We also have several services available to us for extra practice: {\color{blue}\href{https://academicsupport.virginia.edu/math1210}{P2L}}, {\color{blue}\href{https://academicsupport.virginia.edu/mclc-math-collaborative-learning-center}{MCLC}}, {\color{blue}\href{https://academicsupport.virginia.edu/pac}{PAC}}, office hours, etc. Be sure you are incorporating them into your plan! And, yes, you can also incorporate breaks into your study plan.

% \paragraph*{Monday, 25 November}\,

%     \vfill

% \paragraph*{Tuesday, 26 November}\,

%     \vfill

% \paragraph*{Wednesday, 27 November}\,

%     \vfill

% \paragraph*{Thursday, 28 November}\,

%     \vfill

% \paragraph*{Friday, 29 November}\,

%     \vfill

% %%%%%%%%%
% \newpage%
% %%%%%%%%%

% \paragraph*{Saturday, 30 November}\,

%     \vfill

% \paragraph*{Sunday, 1 December}\,

%     \vfill

% \paragraph*{Monday, 2 December}\,

%     \vfill

% \paragraph*{Tuesday, 3 December} (assessment 4 @ 7 PM)\,

%     \vfill

% \paragraph*{Wednesday, 4 December}\,

%     \vfill

% \paragraph*{Thursday, 5 December}\,

%     \vfill

% %%%%%%%%%
% \newpage%
% %%%%%%%%%

% \paragraph*{Friday, 6 December}\,

%     \vfill

% \paragraph*{Saturday, 7 December}\,

%     \vfill

% \paragraph*{Sunday, 8 December}\,

%     \vfill

% \paragraph*{Monday, 9 December}\,

%     \vfill

% \paragraph*{Tuesday, 10 December}\,

%     \vfill

% \paragraph*{Wednesday, 11 December} (final assessment @ 7 PM)\,

%     \vfill

%%%%%%%%%
\newpage%
%%%%%%%%%
    
\item We've talked a lot about mistakes and planning for improvement, but we haven't spent a lot of time considering our successes yet. {\bf Use this space to celebrate your successes!} What are some examples of mathematics you feel like you have learned more about so far this semester? How has your mathematical knowledge grown? Don't be modest - you've earned a chance to talk about yourself positively! :)

    \vfill

\item I have a healthy, productive attitude about learning mathematics. (bubble one) \,\hfill $\bigcirc$ agree \hfill $\bigcirc$ disagree \hfill\,\\[0.1in]
%
Please write anything else you'd like to include in your reflection.

    \vfill

\item The MATH 1190 \& 1210 instruction team wants you to know {\bf we believe in your ability to do great things with enough time and intentional effort}, and {\bf we want to support you in growing throughout this course}. We care, so don't hesitate to reach out if you have any questions or need any help!\\ 

{\bf Please acknowledge this message by signing your name below.}\\[0.5in]
    
\end{enumerate}

\end{document}